% \begin{savequote}
% \includegraphics[width=6cm]{./images/gabriel-sebastian-aass/cape_town_paramedic.jpg}
% 
% {\sffamily\tiny Bild: Gabriel}
% \end{savequote}
\def\ChapterCode{ALS}
\chapter
{Erweiterte Reanimation}
\label{chp:ALS}

\ChapterIntro
{%
}
{%
\begin{description}
  \item[Maintainer:] Sebastian Gabriel
  \item[Autoren:] Diverse
  \item[Reviewer:] Standard-Reviewprozess
  \item[Version:] {\DocVersion} ({\DocVersionInfo})
  \item[SHA1:] {\DocGitCommit}
\end{description}

{\TxTeilkapitelAASS}
}


\clearpage % TEMP temp clearpage
\I{Advanced Life Support} (\I{ALS}) bezeichnet eine Fülle
von erweiterten Maßnahmen zur Therapie des
Herz-Kreislaufstillstands, bzw. im weiteren Sinne von
ähnlichen lebensbedrohlichen Krankheitsbildern. Sie umfassen
unter anderem das Setzen venöser Zugänge, die medikamentöse
Therapie, bestmögliche Atemwegssicherung und Defibrillation.
\Abk{ALS}-Maßnahmen können 
u.\,a. von Sanitätern (ab der Qualifikation \Abk{NKV}) oder
(Not-)Ärzten durchgeführt werden. Der jeweilige Umfang der Maßnahmen
richtet sich nach dem Ausbildungsstand, der eigenen
Erfahrung und dem vorhandenen Material. Sind die
Voraussetzungen für eine Reanimation mittels \Abk{ALS}
nicht gegeben, so muss eine Basisreanimation (wenn
vorhanden mittels \Abk{AED}) durchgeführt werden und weitere
Hilfe angefordert werden. 

\section{Leitlinien}

\Layout{0}{ERC}
{%
\label{topic:erc}
\index{Fachgesellschaft!ERC}
Eine wichtige Fachgesellschaft, welche sich dem Thema der
Reanimation und damit verwandter Gebiete widmet, ist das
\T[ERC|see{European Resuscitation Council}]{ERC}
(\I{European Resuscitation Council}). Bekannt ist das
\Abk{ERC} vor allem durch seine Leitlinien
(\I{Guidelines}\footnote{\Lang{engl.} \emph{Guideline}:
Leit-, Richtlinie, Empfehlung}) zur Wiederbelebung, welche
alle fünf Jahre aktualisiert werden. Die Guidelines
definieren \E{nicht} den einzig richtigen Weg wie eine
Reanimation ablaufen muss. Vielmehr sind sie als weithin
akzeptierte Ansicht zu
verstehen, wie sie sicher und effektiv durchgeführt werden
kann. \cite[p.1219]{Erc2010Sec1}

Jeder neuen Auflage geht ein aufwendiger Forschungs- und
Diskussionsprozess voraus: Es werden hunderte Studien mit
vielen verschiedenen Fragestellungen an vielen
Forschungszentren auf der ganzen Welt durchgeführt, die
Ergebnisse ausgewertet und anschließend diskutiert.
Die derzeit aktuelle Version
der ERC-Guidelines kam im Jahr \EE{2010} heraus
(\I{ERC 2101}), die nächste Aktualisierung ist im Jahr
2015 zu erwarten. Derzeit wird in vielen Bereichen noch nach
der Version aus dem Jahr 2005 gearbeitet.



}{%
\begin{itemize}
   \item Leitlinien zur Reanimation
   \item Aktualisierung alle 5 Jahre
   \item Aktuelle Version: 2010
\end{itemize}

}{}{}{}





\Layout{2}{Neu oder alt?}
{%
Da alle fünf Jahre eine neue Version der Richtlinien
erscheint, kommt es gerade in der
Übergangszeit\footnote{Als Übergangszeit versteht man den
Zeitraum zwischen Veröffentlichung einer neuen Version der
Leitlinie und Implementierung durch die jeweilige
Organisation.} regelmäßig zu Verwirrungen, welche Versionen
oder Varianten in der Praxis angewendet werden sollen.

\EE{Grundsätzlich muss das gesamte Personal einer Organisation
\E{das von der eigenen Organisation empfohlene} oder
angepasste Schema kompetent im Team anwenden können.} Das
behandelnde Fachpersonal kann jedoch davon im Rahmen der
\E{Eigenverantwortung} abweichen, sofern die dazu
notwendigen Voraussetzungen gegeben sind. Dabei ist ganz
besonders zu beachten, dass das \E{gesamte} Team mit der
anderen oder neueren Version vertraut und die entsprechende
Routine besitzen muss. Darüber hinaus müssen auch
z.\,B. technische Gegebenheiten bedacht werden
(Programmierung der AEDs, etc.). I.\,d.\,R. trifft der
Teamleiter als Letztverantworlicher die
Entscheidung, er muss daher auch genau über die jeweilige
Qualifikation und Routine seiner Teammitglieder Bescheid
wissen. Ist die Durchführung einer anderen Variante als die
von der Organisation empfohlene aufgrund der Ausbildung, des
Trainings und der Routine von Team-Mitgliedern oder der
vorhandenen Ausstattung fragwürdig, so ist im Zweifel im Sinne eines
reibungslosen Ablaufes des empfohlene Schema zu verwenden.
\E{Die Publikation einer neuen Ausgabe der Guidelines
bedeutet nicht, das die bisherigen Methoden unsicher oder
ineffektiv sind \cite[p.1219]{Erc2010Sec1}.}

% \enquote{The guidelines [..] do not define the only way that
% resuscitation can be delivered; they merely represent a
% widely accepted view of how resuscitation should be
% undertaken both safely and effectively. The publication of
% new and revised treatment recommendations does not imply
% that current clinical care is either unsafe or ineffective.}

}{%
\begin{itemize}
   \item
\end{itemize}

}{}{}{}

\bigskip

% \EE{Sonderfall: VF/VT unter bestimmten Bedingungen}:
% Dies betrifft
% Patienten mit Kammerflimmern oder ventrikulärer
% Tachykardie \Z{während einer Herzkatheteruntersuchung oder
% früh-postoperativ nach einem herzchirugischen Eingriff bei
% einem beobachteten Kreislaufstillstand, wenn der Patient
% bereits an einem manuellen Defibrillator angeschlossen
% ist: Hier dürfen 3 Schocks hintereinander gegeben
% werden} Unter Bedachtnahme des Abschnitts bezüglich der
% Elektrotherapie gilt dies auch für Patienten, welche einen
% beobachteten Kreislaufstillstand erleiden und zu dem
% Zeitpunkt bereits an einem Defibrillator angeschlossen
% sind.

\clearpage
\section{Säulen der Reanimation}

\subsection{Basisreanimation}

Das unter \CO{[EHI]} (\prettyref{chp:EHI}) und \CO{[BLS]}
(\prettyref{chp:BLS}) Gesagte gilt auch beim Advanced
Life Support. Auch hier muss besonders darauf geachtet
werden, so wenig wie möglich zu pausieren und
Herzdruckmassagen mit möglichst optimaler Qualität
durchzuführen. Unterbrechungen sind nur statthaft, um
spezielle Interventionen zu ermöglichen.

\Layout{2}{Neuerungen 2010}{%
Die Empfehlung des
\enquote*{Pflicht-Zyklus vor AED} beim unbeobachteten
Kreislaufstillstand ist gestrichen
worden, kann aber weiterhin angewendet werden.
Während der Ladephase des
Defibrillators soll die Herzdruckmassage (\I{HDM})
fortgeführt werden um die HDM-Pausen zu reduzieren,
d.\,h. nach der Analyse wird weiter massiert, bis der Defi bereit zur
Schockabgabe ist.  Bei
vielen AED-Geräten kann es jedoch zu
\EE{Problemen} kommen, sie interpretieren die
reanimationsbedingten Störungen der EKG-Ableitung als
Kontaktverlust bzw. Bewegungsartefakt und brechen den
Ladevorgang ab.\footnote{Davon betroffen sind derzeit z.\,B. die
Geräte der Schiller Fred{\texttrademark}-Serie, der
Defigard{\texttrademark} 1002 und der
LifePak{\texttrademark} 500. Bezüglich anderer Geräte liegen
der Redaktion derzeit noch keine Informationen vor. Es muss
jedoch davon ausgegangen werden, dass auch andere Modelle
ein derartiges Verhalten zeigen können. \E{(Angaben ohne
Gewähr. Stand: 2011-01-19)}} Dadurch wird die Schockabgabe
verhindert. Daher kann diese Methode nicht generell
empfohlen werden, sondern darf nur durchgeführt werden, wenn
der Hersteller des Gerätes dies freigegeben hat.
\Commentary[HDM während Defi lädt / Doppelte Warnung]{Eine
Warnung betreffend der HDM während der Ladephase des Defis
findet sich bereits im \CO{[BLS]}-Teil
(\prettyref{chp:BLS}). Aufgrund der Tragweite wurde sie im
\CO{[ALS]}-Teil praktisch
wortwörtlich wiederholt.}

\Cave{Manche Defibrillatoren brechen den Ladevorgang ab,
wenn währenddessen eine Herzdruckmassage durchgeführt wird.}

\Fazit{Während des Ladevorganges des Defibrillators darf nur
bei Verwendung von \E{manuellen} Geräten eine
Herzdruckmassage durchgeführt werden.}
}{}{}{}{}


\subsection{Elektrotherapie}

Vgl. \CO{[BLS]} (\prettyref{chp:BLS}).

% \TakeNotesLarge

\subsection{Medikamente}
\index{Adrenalin!Reanimation}%
\index{Amiodaron!Reanimation}%
\index{Atropin!Reanimation}%
Die Gabe von Medikamenten ist eine Routinemaßnahme während
der Reanimation. Typischerweise erfolgt die Medikamentengabe
über eine periphere Venenverweilkanüle
(Venflon{\texttrademark}, Braunüle{\texttrademark}).
Sinvoller weise wird standardmäßig auch eine Infusion mit
kristalloider Infusionslösung (Kochsalz-,
Ringerlösung, Ringerlaktat) angeschlossen, um eine schnelle
Spülung zu ermöglichen.

Sollte
die rasche Anlage einer solchen nicht gelingen, kann als
Alternative ein intraossärer Zugang mittels spezieller
Hilfsmittel gelegt werden.
Früher war es auch üblich, Medikamente über einen
endotracheal-Tubus zu geben. Dies wird nicht mehr
empfohlen.

% \Layout{2}{Medikamente ERC 2005}{%
%
% \subparagraph{Schockbarer Rhythmus} Bei
% schockbaren Rhythmen (VF/VT) gilt nun:
% \begin{labeling}[:]{Adrenalinm}
%   \item[Adrenalin] 1\Mg* vor dem 3. Schock, wenn
%   die HDM fortgeführt wird, danach alle
%   \VonBis{3}{5}\Min* während wechselnder CPR-Zyklen.
%   \item[Amiodaron] 300\Mg* vor dem 4. Schock.
% \end{labeling}
%
% \subparagraph{Nichtschockbarer Rhythmus} (Asystolie, Pulslose
% Elektrische Aktivität (\I{PEA})): Bei nichtschockbaren
% Rhythmen gilt:
% \begin{labeling}[:]{Adrenalinm}
%   \item[Adrenalin] 1\Mg* sobald Zugang vorhanden ist,
%   danach alle \VonBis{3}{5}\Min*.
%   \item[Atropin] 3\Mg* sobald Zugang vorhanden ist,
%   einmalig.
% \end{labeling}
% }
% {}{}{}{}
%
% \Layout{2}{Medikamente ERC \Ca{2010}}{%
% \subparagraph{Schockbarer Rhythmus} Bei
% schockbaren Rhythmen (VF/VT) gilt nun:
% \begin{labeling}[:]{Adrenalinm}
%   \item[Adrenalin] 1\Mg* \Ca{nach} dem 3. Schock, wenn
%   die HDM fortgeführt wird, danach alle
%   \VonBis{3}{5}\Min* während wechselnder CPR-Zyklen.
%   \item[Amiodaron] 300\Mg* \Ca{nach} dem \Ca{3.} Schock.
% \end{labeling}
%
% \subparagraph{Nichtschockbarer Rhythmus} (Asystolie, PEA): Bei
% nichtschockbaren Rhythmen  gilt nun:
% \begin{labeling}[:]{Adrenalinm}
%   \item[Adrenalin] 1\Mg* sobald Zugang verfügbar, danach alle
%   \VonBis{3}{5}\Min* während wechselnder CPR-Zyklen.
% \end{labeling}
%
% \noindent\EE{Atropin fliegt raus}: Atropin ist nun \E{nicht}
% mehr als Routine-Medikament bei nichtschockbaren Rhythmen
%  empfohlen.
% }
% {}{}{}{}


\begin{table}[H]
\FontTableBase
\begin{WidePar}
\Captionof{table}{Übersicht: Die wichtigsten bei der
Reanimation eingesetzten
Medikamente.}

\begin{tabulary}{\linewidth}{lLLL}
\toprule\addlinespace
\noindent\TabHeading{Wirkstoff}
&
\noindent\TabHeading{Erläuterungen}
&
\noindent\TabHeading{Indikation}
&
\noindent\TabHeading{Spezialitäten}
\\\addlinespace\midrule\addlinespace
\noindent\TabSub{Adrenalin}
&
Vasokonstriktor
&
Schockbarer Rhythmus

Nicht-Schockbarer Rhythmus
&
\noindent\RaggedRight\Z{L-Adrenalin Fresenius{\texttrademark} 2\Mg*\,/\,20\Ml*

L-Adrenalin Fresenius{\texttrademark} 0,5\Mg*\,/\,5\Ml*

Suprarenin{\texttrademark} 1\Mg*\,/\,1\Ml*}
\\\addlinespace
\noindent\TabSub{Amiodaron}
&
Anti\-ar\-rhythmi\-kum

Kl. III
&
Schockbarer Rhythmus
&
\noindent\RaggedRight\Z{Sedacoron{\texttrademark} 150\Mg*\,/\,3\Ml*}
\\\addlinespace
\noindent\TabSub{\Z{Atropin}}
&
\noindent\Z{Anti\-para\-sym\-patho\-mi\-meti\-kum}
&
\noindent\Z{Nicht-Schockbarer Rhythmus (seit 2010 nicht mehr
Standard)}
&
\noindent\RaggedRight\Z{Atropinum sulfuricum Nycomed{\texttrademark} 0,5\Mg*\,/\,1\Ml*}
\\\addlinespace\bottomrule
\end{tabulary}
\end{WidePar}
\end{table}









% \Layout{57}
% {}
% {}
% {}
% {
% \begin{tabularx}{\linewidth}{XXX}
% \TabHeading{Medikament}
% &
% \TabHeading{Schockbar}
% &
% \TabHeading{Nicht-schockbar}
% \\\midrule\addlinespace
% \TabSub{Adrenalin}
% &
% \EE{1\Mg*} \Abk{IV}
%
% vor dem 3., 5., 7.,~{\ldots} Schock
% &
% \EE{1\Mg*} \Abk{IV}
%
% ab Zugang, alle \VonBis{3}{5} Minuten
% \\\addlinespace
% \TabSub{Amiodaron}
% &
% \EE{300\Mg*} \Abk{IV}
%
% einmalig vor \E{4.} Schock
% &
% ---
% \\\addlinespace
% \TabSub{Atropin}
% &
% ---
% &
% \EE{3\Mg*} \Abk{IV}
%
% ab Zugang
% \\\addlinespace\bottomrule\addlinespace
% \end{tabularx}
% }
% {Standard-Reanimationsmedikamente ERC 2005 (leicht vereinfacht)}
% {}


\Layout{57}
{}
{}
{}
{
\begin{tabularx}{\linewidth}{XXX}
\toprule\addlinespace
\noindent\TabHeading{Medikament}
&
\noindent\TabHeading{Schockbar}
&
\noindent\TabHeading{Nicht-schockbar}
\\\addlinespace\midrule\addlinespace
\TabSub{Adrenalin}
&
\EE{1\Mg*} \Abk{IV}

\Ca{nach} dem 3., 5., 7.,~{\ldots} Schock
&
\EE{1\Mg*} \Abk{IV}

ab Zugang, alle \VonBis{3}{5} Minuten
\\\addlinespace
\TabSub{Amiodaron}
&
\EE{300\Mg*} \Abk{IV}

einmalig \Ca{nach} dem \Ca{3.} Schock
&
---
\\\addlinespace\bottomrule
\end{tabularx}
}
{Standard-Reanimationsmedikamente ERC \Ca{2010} (leicht vereinfacht)}
{}

\subsection{Airwaymanagement}
Es gibt nicht genügend Beweise
für oder gegen eine spezielle Technik zum
\I{Airwaymanagement} beim Kreislaufstillstand.

\Layout{2}{Endotracheale Intubation}
{Die frühe endotracheale Intubation (Einführen eines Tubuses
in die Luftröhre) verliert an Stellenwert. Sie soll nur
angewendet werden, wenn der Anwender hochqualifiziert
(\Speech{\enquote{high level of skill {\ldots}}}
\cite{Erc2010Sec4}) und sich sehr sicher ist
(\Speech{\enquote{{\ldots} and confidence}}
\cite{Erc2010Sec4}), bzw. regelmäßige, laufende Erfahrung
mit der Technik hat (\Speech{\enquote{and has regular,
ongoing experience with the technique.}}
\cite{Erc2010Sec4}).
Ohne ausreichendes Training und Erfahrung wird die
Komplikationsrate als unakzeptabel hoch angesehen. Das
Aufschieben der endotrachealen Intubation bis zur Wiederkehr
des spontanen Kreislaufs, wird in den ERC 2010-Leitlinien
als Alternative ausdrücklich genannt.
\cite{Nolan2001,
Nolan2008,
Wang2006,
Wang2006a,
Stiell2004,%There was no improvement in the rate of survival
% with the use of advanced life support in any subgroup.
Lecky2008,
Erc2010Sec4}
\Commentary[Airwaymanagement / Frühe Endotracheale Intubation (fETI)]
{Die endotracheale Intubation (ETI) ist grundsätzlich eine
sehr anspruchsvolle Tätigkeit, besonders im
außer-klinischen Bereich; es gibt viele Fehlerquellen
(Intubation des Ösophagus, Intubation eines Bronchus,
iatrogene Verletzung beim Intubationsvorgang, \ldots).
Zahlreiche Studien haben eine hohe Inzidenz von
Fehlintubationen beschrieben.

Nolan 2001 \cite{Nolan2001}
gibt eine umfassende
Übersicht bezüglich Erfolgsraten und Trainingsstrategien.
Schon damals waren die dargestellten Ergebnisse nicht
zufriedenstellend, gleichzeitig aber durchwachsen. Im
weiteren zeitlichen Verlauf wurde die ETI immer weiter
hinterfragt, z.\,B.
\cite{Wang2006,Wang2006a,Wang2008,Stiell2004,Stiell2008}.

In 2008 fassten Nolan und Soar \cite{Nolan2008} zusammen,
dass der Nutzen der ETI bei der präklinischen Reanimation
nicht durch Studien gezeigt werden konnte, und dass einige
sogar das Gegenteil erkennen ließen.
(\Speech{\enquote{Prehospital studies fail to show any
benefit from tracheal intubation during cardiopulmonary
resuscitation and many show harm.}}~\cite{Nolan2008}) Ein
ebenso 2008 veröffentlichtes Cochrane-Review von Lecky
et~al. \cite{Lecky2008} kam zu der gleichen Schlussfolgerung
mit der Betonung, dass entsprechende, qualitative
hochwertige Studien fehlen.
Dass die Probleme der ETI nicht auf nicht-ärztliches
Personal beschränkt ist, zeigten Timmermann et.~al. 2007
\cite{Timmermann2007}. In dieser in Deutschland
durchgeführten Studie wurde eine hohe Rate an
Fehlintubationen (ösophageal, bronchial) durch Notärzte
festgestellt.

Vor diesem Hintergrund bleibt festzuhalten, dass eine
unterlassene Intubation nicht mit einer unterlassenen
Hilfeleistung gleichzusetzen ist.}
}{}{}{}{}{}

\subparagraph{Kontrolle der Tubuslage}
Das korrekte Platzieren des Tubus ist einer der kritischten
Punkte bei der endotrachealen Intubation. Typische Fehler
sind die Intubation der Speiseröhre (\E{ösophageale}
Intubation) oder die Intubation eines Hauptbronchus
(\E{bronchiale} Intubation, wenn der Tubus zu weit
vorgeschoben wird; es wird dann nur \E{ein} Lungenflügel belüftet).

Grundsätzlich gibt es präklinisch drei Möglichkeiten, die
korrekte Lage des Tubus zu überprüfen:
\begin{enumerate}
  \item Visuelle Darstellung mittels Laryngoskop
  \item Auskultation von Magen und Lungenflügel
  \item Messung der \ce{CO2}-Abatmung mittels Kapnograhie
\end{enumerate}

 Grundsätzlich sollen immer \EE{alle drei
Möglichkeiten} angewendet werden.
Es wird jedoch zunehmend ein größerer Wert auf die
\EE{Kapnographie} gelegt, um die Tubuslage zu kontrollieren,
die CPR-Qualität abzuschätzen und eine frühzeitige Erkennung
der Rückkehr des spontanen Kreislaufs zu ermöglichen.
Besonders geeignet sind dazu Geräte, welche die
\ce{CO2}-Abatmung als Kurve darstellen. Wenn nach 6
Beatmungen \ce{CO2} nachweisbar ist, kann man davon
ausgehen, dass der Tubus in den unteren Atemwege zu liegen
gekommen ist. Die Kapnographie kann jedoch \E{keine} Aussage
darüber treffen, ob die \E{Trachea oder ein einzelner
Bronchus} intubiert wurde.

Wenn kein Kapnograph mit Kurvendarstellung vorhanden ist,
wird in den ERC~2010-Leitlinien die bevorzugte Verwendung
supraglottischer Hilfsmittel
angeregt.\ZFootnote{\Speech{\enquote{In the absence of a
waveform capnograph it may be preferable to use a
supraglottic airway device when advanced airway management
is indicated.}}~\cite[p.1322]{Erc2010Sec4}.
\Commentary[Airwaymanagement / Kapnographie]{Zur
Problematik der Verfügbarkeit der Kapnographie siehe
auch 
\cite{Genzwuerker2007}
(\fullcite{Genzwuerker2007}) 
\cite{Timmermann2007}
und darin enthaltene
Referenzen.
}
}




\Layout{2}{Supraglottische Atemwegshilfsmittel}
{In den letzten Jahren gewinnen die
supraglottischen\footnote{\Lang{lat.} supra: oberhalb, über;
\Lang{Term.} glottis: Der aus den Stimmfalten bestehende
Teil des Kehlkopfinnenraums.} Hilfsmittel aufgrund ihrer
verhältnismäßig einfachen Anwendbarkeit immer mehr an
Bedeutung. Die ERC-Leitlinien 2010 sehen mittlerweile die
supraglottischen Atemwegshilfen bei der Reanimation als
weitgehend gleichwertig zur klassischen endotrachealen
Intubation an
\cite[p.1319,1321]{Erc2010Sec4}.\footnote{Supraglottische
Methoden zur Atemwegssicherung werden oft als
\enquote*{alternativer Atemweg} bezeichnet, da sie oft eine
Alternative zur endotrachealen Intubation (\I{ETI})
angesehen werden. Die ETI gilt weithin als Goldstandard der
Atemwegssicherung, dennoch hat sie einige gewichtige
Nachteile und ist in vielen Situationen \E{nicht die Methode
der Wahl}. }

Bei den supraglottischen Atemwegshilfen wird der Schlauch
nicht in die Luftröhre eingeführt, sondern endet im Rachen
beziehungsweise oberhalb des Kehldeckels (supraglottisch).
Es gibt verschiedene Produkte, welche dieses Ziel auf
unterschiedliche Weise erreichen. Dazu gehören der
Larynxtubus{\texttrademark}, die Larynxmaske und der Combitubus.

 }{}{}{}{}{}


\subsubsection{Larynxtubus{\texttrademark}}

\Layout{2}{{\TxBeschreibung} Larynxtubus{\texttrademark}}
{\label{topic:larynxtubus}
Der \T{Larynxtubus{\texttrademark}} wurde von der Firma VBM
Medizintechnik entwickelt und zählt derzeit zu den
populärsten Produkten zur supraglottischen
Atemwegssicherung. Im Gegensatz zur endotrachealen
Intubation kann er blind und ohne Hilfsmittel eingeführt
werden, das Ende kommt bei korrekter Anwendung in der
Speiseröhre zu liegen. Er besitzt \E{ein}
Lumen\footnote{Als \I{Lumen} bezeichnet man denn
Innendurchmesser von Röhren, Kanülen und Kathetern.} und
zwei Cuffs\footnote{\I{Cuff}:
ein mit Luft gefüllter Ballon.}, manche Modelle (LT-S) verfügen über ein zweites Lumen, über
welches eine Magensonde zur Ableitung des Mageninhaltes
gelegt werden kann. Der obere Cuff kommt nach dem Einführen
im Rachen zu liegen, während der untere am Ende die
Speiseröhre abdichtet.
Dazwischen liegt die Öffnung des Lumens, nach dem Einführen
in Höhe des Kehlkopfes. Da ober- und unterhalb der
Abzweigung zur Luftröhre abgedichtet ist, kann die Luft über
den Kehlkopf und den Kehldeckel in die Luftröhre strömen und
die Gefahr der Aspiration von Mageninhalt kann deutlich
reduziert werden. Es sind verschiedene Größen je nach
Körperlänge erhältlich.
\cite{Asai2005} % The laryngeal tube.

Im direkten Vergleich mit der Beatmung mittels
Beatmungsmaske
sind bei einer korrekten Anwendung des Larynxtubus{\texttrademark} folgende
Vorteile gegeben:

\begin{itemize}
  \item Durch den geblockten Ösophagus wird ein Überblähen
  des Magens sowie die Aspiration von Mageninhalt
  verhindert.
  \item Die Anwendung erfordert keine höhere manuelle
  Geschicklichkeit als die Anwendung der Beatmungsmaske und
  -beutel.
  \item Zusätzliche Hilfsmittel zur Anwendung sind in der
  Regel nicht erforderlich, wodurch die Gefahr von
  Verletzungen durch solche reduziert ist.
\end{itemize}

Der Larynxtubus{\texttrademark} ist während der Reanimation
eine geeignete Alternative zur Maskenbeatmung und zur
endotrachealen Intubation und kann auch durch Angehörige der
nichtärztlichen Gesundheitsberufe im Rahmen der
lebensrettenden Sofortmaßnahmen angewendet werden.
\Commentary[Airwaymanagement / Larynxtubus /
Stellungsnahme]{Gemäß einer Stellungsnahme des BM für
Gesundheit ist eine Anwendung des Larynxtubus für Angehörige
der nichtärztlichen Gesundheitsberufe im Rahmen der
lebensrettenden Sofortmaßnahmen möglich, bedarf aber
spezieller Maßnahmen im Rahmen der Aus-, Fort- und
Weiterbildung. Vgl. Schreiben vom 10.12.2010, Geschäftszahl
BMG-92250/0080-II/A/2010)}
Eine Anwendung setzt jedoch entsprechende
Kenntnisse und Fertigkeiten voraus. Diese müssen im Rahmen
von speziellen Fortbildungen erworben werden; d.\,h. die
Anwendung setzt eine erfolgte Einschulung voraus.
\cite{BMG-92250_0080-II_A_2010_2010-12-15}
}{}{}{}{}{}

\begin{PictureSeriesWide}{}{Übersicht: Lage verschiedener
Hilfsmittel zur Atemwegssicherung
\SourcePicture{Hirtler}\label{fig:uebersicht-lage-atemwegshilfen}}
\PictureSeriesRowWide{3}
{./images/hirtler-aass/PositionGuedel-color}
{Position des Guedel-Tubus im Querschnitt}
{./images/hirtler-aass/LageLarynxTubus}
{Im Vergleich: Lage eines Larynx-Tubus \ldots}
{./images/hirtler-aass/LageTubus-color}
{\ldots\ und eines endotrachealen Tubus.}
{empty}
{}
\end{PictureSeriesWide}


\Layout{47}
{}
{}
{}
{\begin{tabulary}{\linewidth}{LLLCL}
\addlinespace\toprule\addlinespace
\noindent\TabHeading{Größe}
&
\noindent\TabHeading{Farbe}
&
\noindent\TabHeading{Patientengruppe}
&
\noindent\TabHeading{Größe}
&
\noindent%\TabHeading{Cuff \Ml}
\\\addlinespace\midrule\addlinespace
\TabSub{0}	& \TabSub{Transparent}	& Neugeborene			& \textless\,5\Kg*
\\
\TabSub{1}	& \TabSub{Weiß}			& Kleinkinder			& \VonBis{5}{12}\Kg*
\\
\TabSub{2}	& \TabSub{Grün}			& Kinder				& \VonBis{12}{25}\Kg*
\\
\TabSub{3}	& \TabSub{Gelb}			& Erwachsene (klein)	& \textless\,155\CM*
\\
\TabSub{4}	& \TabSub{Rot}			& Erwachsene (mittel)	& \VonBis{155}{180}\CM*
\\
\TabSub{5}	& \TabSub{Violett}		& Erwachsene (groß)		& \textgreater\,180\CM*
\\\addlinespace\bottomrule\addlinespace
\end{tabulary}}
{Larynxtubus: Größen}
{}

\begin{PictureSeriesWide}{}{Bilderserie: Larynxtubus{\texttrademark}
\SourcePicture{Ch. Pallinger}}
\PictureSeriesRowWide{3}
{./images/pallinger-christoph-aass/Larynxtubus_32890_textwidth}
{Larynxtubus{\texttrademark} mit Blockerspritze \ldots}
{./images/pallinger-christoph-aass/Larynxtubus_32885_textwidth}
{{\ldots} und geblocktem Cuff.}
{./images/pallinger-christoph-aass/Larynxtuben_Set_33025_textwidth}
{Larynxtuben gibt es in verschiedenen Größen.}
{empty}
{}
\end{PictureSeriesWide}



\Layout{2}
{Material}
{
\begin{multicols}{2}
\begin{compactenum}
  \item Beatmungsbeutel mit Maske, \ce{O2}-Line und Reservoir
  \item Absaugung
  \item Fixierung
  \item Passender Larynxtubus mit Blockerspritze
  \item Stethoskop
  \item Gleitmittel
  \item Bakterienfilter mit \I{Gänsegurgel}
  \item Magensonde
  \item Evtl. Kapnometrie
\end{compactenum}
\end{multicols}
}
{}
{}
{}
{}



\Layout{2}
{Durchführung}
{Der
Larynxtubus wird ohne Hilfsmittel
über den Mund des Patienten eingeführt, wobei der Kopf
in \T{Schnüffelstellung} (\I{Jackson-Position}, Kopf
  leicht erhöht (umgedrehte Nierentasse), nicht überstreckt,
  nicht gebeugt) gelagert wird. Zusätzliche Hilfsmittel
 sind in der Regel nicht erforderlich.

Das Einführen erfolgt vorsichtig, ohne Gewalt, bei
Widerstand muss gestoppt werden. Es sollen nicht mehr als 2
Versuche gemacht werde. Ein
Versuch darf maximal 30 Sekunden dauern, dazwischen muss der
Patient mit Beutel-Masken-Beatmung oxygeniert werden
Eine laufende CPR darf durch die Intubation mit dem
Larynxtubus nicht unterbrochen werden.

\begin{WidePar}
\begin{multicols}{3}
\begin{enumerate}\CorrectionListsBox
  \item \EE{Notarzt} alarmieren
  \item \EE{Lagerung}: Rückenlage, Kopf in
  \EE{Schnüffelstellung} (Jackson Position, Kopf
  leicht erhöht (umgedrehte Nierentasse), nicht überstreckt,
  nicht gebeugt)
  \item Mund öffnen mit dem
  \enquote*{Kreuzgriff} (Daumen und Zeigefinger der linken
  Hand, der Daumen am Unterkiefer, der Zeigefinger am
  Oberkiefer (bzw. Zähne) öffnen den Mund. Achtung bei
  kariösen oder wackeligen Zähnen!)
  \item \EE{Einführen}: LT mit der Spitze
  entlang des harten Gaumens einführen, ev. Unterkiefer
  anheben
  \item Zahnreihe innerhalb der Markierungen
  \item
  \EE{Cuffen}, Volumen nach Farbcodierung
  \item Manuelles
  \EE{Fixieren}
  \item \EE{Beatmen}
  \item \EE{Erfolgskontrolle}: sehen,
  hören, \ce{CO2}-Messung, Sp\ce{O2} steigt, gleichbleibend
  \begin{enumerate}
    \item Sehen: Thoraxexkursionen
    \item Hören: Auskultation
    Kein Geräusch über dem
    Magen, Atemgeräusche über beiden Lungenflügeln.

    Achtung: Nicht seitengleiches Atemgeräusch beim
    Beatmen über den LT heißt: Verlegung eines Hauptbronchus
    (Fremdkörper) oder Pneumothorax, ist jedoch kein
    Hinweis für einen zu tief oder falsch liegenden Tubus
    (wie bei der Auskultation nach der endotrachealen
    Intubation))!
    \item Messung: sofortige 2-stellige et\ce{O2}-Messung,
    Sp\ce{O2} steigend oder gleichbleibend, Verfärbung des EasyCap
    II
  \end{enumerate}
  \item Wenn Beatmung gut möglich: \EE{Fixierung}
  \item Vorsichtiges Einführen einer \EE{Magensonde} und Entlasten des
  Magens (Magensonde erst einführen, wenn der Patient
  verlässlich beatmet ist und der LT fixiert ist (kein Muss!)
  \item Maximal 2 Versuche, max. 30\Sec* pro Versuch
  \item Absaugung von Sekret aus LT: Sekret aus dem LT  darf
  abgesaugt werden. Vorher die Absaugtiefe mit dem Abstand
  Mundwinkel-Ohrläppchen bestimmen
\end{enumerate}
\end{multicols}
\end{WidePar}




\Fazit{Nur durch alle angeführten Punkte kann eine Fehllage ausgeschlossen werden}}
{}
{}
{}
{}

\Layout{2}
{Tipps, Tricks, Troubleshooting}
{
\begin{itemize}
  \item Beim Einführen des LT auf die Zähne achten, diese
  können den Cuff beschädigen
  \item Durch Hochziehen des Unterkiefers wird der
  Zungengrund angehoben, das kann Einführen des LT
  erleichtern
  \item Eventuell die Zunge mit den Fingern der li. Hand
  etwas zur Seite schieben, um das Einführen zu erleichtern
  \item Möglicherweise Neupositionieren des Kopfes während
  des Einführens
\end{itemize}


}
{}
{}
{}
{}

\Layout{2}
{Komplikationen}
{Grundsätzlich ist der Larynxtubus vergleichsweise sicher in
der Anwendung. Komplikationen umfassen eine Undichtigkeit
(Leakage), Schädigung der Rachenschlimhaut, evtl. mit
Blutung, und eine Überblähung des Magens und eine Ruptur
der Speiseröhre. Der Larynxtubus bietet zwar einen guten
Aspirationsschutz, dieser ist jedoch nicht komplett
zuverlässig.}
{}
{}
{}
{}

\Layout{2}
{Beatmung über den Larynxtubus}
{
\begin{itemize}
  \item 100\PC* Sauerstoff
  \item Reanimation weiter mit 30:2 fortsetzen
  \item Frequenz \VonBis{10}{15}\PerMin* (bei ROSC)
  \item Atemzugvolumen ca. \VonBis{350}{500}\Ml*, Thoraxexkursion kontrollieren

\end{itemize}

\Customized{
\begin{description}
  \item[\CustAsboe]
  Die Beatmung über den Larynxtubus darf für NFS ohne NKI
  nur über den Beatmungsbeutel erfolgen.
\end{description}
}
}
{}
{}
{}
{}







\Customized{
\begin{description}
  \item[\CustAsboe] Die Bundesschulung und die Chefärzte des
  Bundesverbandes haben den Larynxtubus zur Anwendung unter
  bestimmten Voraussetzungen für den Rettungsdienst freigegeben.

  Voraussetzungen für die Anwendung des
  Larynxtubus ist, dass der Anwender
  \EE{Notfallsanitäter} ist oder eine höhere Qualifikationen
  besitzt, und eine \EE{entsprechende
  Ausbildung}
  absolviert hat. Diese
  zusätzliche Kompetenz muss auch bei der
  erforderlichen Rezertifizierung ausreichend
  bei einem Fallbeispiel überprüft werden.
  Notfallsanitäter mit Notfallkompetenz
  Beatmung und Intubation benötigen diese
  zusätzliche Ausbildung nicht.

  Die Anwendung eines Larynxtubus \EE{durch
  Rettungssanitäter ist vorerst nicht gestattet}.
  \cite{ASBOeRundschreiben2010Nr11}
\end{description}
}




\subsection{Reversible Ursachen}

Es gibt acht klassische Ursachen eines
Herz-Kreislaufstillstandes, welche grundsätzlich behebbar
(reversibel) sind:

\begin{multicols}{2}
\TableBaseModification
\begin{itemize}
  \item \EE{\Ca{H}ypoxie}

  (z.\,B. Atemwegsverlegung),
  \item \EE{\Ca{H}ypovolämie}

  (z.\,B. Blutung durch Trauma,   GI-Blutung, Aneurysmaruptur),
  \item \EE{\Ca{H}ypo-/Hyperkaliämie}

  (z.\,B. metabolische Störung),
  \item \EE{\Ca{H}ypothermie}

  (z.\,B. Ertrinkungsunfall)
  \item \EE{\Ca{H}erzbeuteltamponade}

  (z.\,B. Thoraxtrauma durch Stich oder Schuss)
  \item \EE{\Ca{I}ntoxikation}

  (z.\,B. Opiate, trizyklische Antidepressiva)
  \item \EE{\Ca{T}hromboembolie}

  (Lungenembolie, Herzinfarkt)
  \item \EE{\Ca{S}pannungspneumothorax}
\end{itemize}
\end{multicols}


\noindent \Sign{→} \T{4 H + HITS}

\bigskip

Leider lassen sich nicht alle dieser reversiblen Ursachen
präklinisch erkennen oder beheben, dennoch muss dies -- wenn
möglich -- versucht werden.


\section{Versorgung nach Reanimation}

Wird bei der Reanimation der Kreislauf wiederhergestellt
(\T{ROSC}, \I{Resumption Of Spontaneous Circulation}) ist
die weitere Versorgung wichtig. Grundsätzlich gelten auch
hier sinngemäß die {\TxMassVit}. Die Sauerstoffgabe hat
einen Sp\ce{O2}-Zielwert von \VonBis{94}{98}\PC*.
Dem Post-Cardiac-Arrest-Syndrome sowie dessen Behandlung
wird vermehrt Beachtung geschenkt.

Strukturierte post-Resuscitation-Protokolle erhöhen die
Überlebenswahrscheinlichtkeit und werden zunehmend
implementiert. Besonders aktuell diesbezüglich ist die
\T{Therapeutische Hypothermie}, bei welcher die
Körpertemperatur des Patienten kontrolliert reduziert wird,
um die körpereigene Bildung von schädlichen Stoffen nach
einer Reanimation zu mildern.



\section{Abbruch der Reanimation}

In ca. \VonBis{70}{95}\PC* der Fälle ist der
Reanimationsversuch erfolglos und der Patient verstirbt
\cite{StriebelAnaesthesie2-Band2,Erc2005DeSec8}.
Grundsätzlich wird eine Reanimation fortgesetzt, solange ein
Kammerflimmern oder ein ähnlicher Rhythmus vorliegt. Es ist
weitgehend akzeptiert, eine Reanimation abzubrechen, wenn
trotz aller ALS-Maßnahmen und Fehlen von reversiblen
Ursachen für 20 Minuten eine Asystolie besteht
\cite{Erc2010Sec10,Bonnin1993,StriebelAnaesthesie2-Band2}.
Gewisse Bedingungen, etwa eine Hypothermie zum Zeitpunkt
des Kreislaufstillstands, steigern die Chancen der Wieder-
herstellung ohne neurologische Schäden,
sodass normale Prognosekriterien (z.\,B.
eine über 20\Min* andauernde Asystolie)
nicht anwendbar sind. \autocite{Erc2005DeSec8}

Abgesehen von den rein medizinischen Faktoren,
welche vor allem den zu erwartenden Reanimationserfolg im
Blick haben, gibt es auch noch andere Umstände, welche
Einfluss auf die Entscheidung haben. Besonders
berücksichtigenswert ist der mutmaßliche oder vorab
geäußerte Patientenwille, z.\,B. im Rahmen einer
verbindlichen oder beachtlichen Patientenverfügung.

Die Abbruchsentscheidung ist grundsätzlich \EE{von einem
Arzt} zu treffen (sofern keine dringenden Gründe wie
unsichere Umgebung oder Erschöpfung vorliegen).

\opt{STATUS-EXPERIMENTAL}{

}%STATUS-EXPERIMENTAL

\clearpage
\begin{figure}
\begin{WidePar}
\begin{center}
\Captionof{figure}{Algorithmus Advanced Live Support (\Abk{ALS})}

\includegraphics[scale=1]{./images/emhofer-josef-aass/ALS}
\end{center}
\end{WidePar}
\end{figure}

\nocite{Uray2008}


\nocite{Erc2010Sec1}
\nocite{Erc2010Sec2}
\nocite{Erc2010Sec3}
\nocite{Erc2010Sec4}
\nocite{Erc2010Sec5}
\nocite{Erc2010Sec6}
\nocite{Erc2010Sec7}
\nocite{Erc2010Sec8}
\nocite{Erc2010Sec9}
\nocite{Erc2010Sec10}
\ChapterEnd
